\begin{trapbox}
\begin{itemize}
    \item \textbf{混淆"否命题"与"命题的否定"}：否命题是对条件和结论都否定，而命题的否定（\(\lnot(p\to q)\)）只否定结论（且保持条件不变），二者完全不同。例如原命题"若 \(x>0\) 则 \(x>1\)"的否命题是"若 \(x\le0\) 则 \(x\le1\)"，而命题的否定是"存在 \(x>0\) 使得 \(x\le1\)"。
    \item \textbf{误认为逆命题与否命题等价于原命题}：原命题与逆否命题等价，逆命题与否命题等价，但原命题与逆命题不一定等价。解题时不可随意互换。
    \item \textbf{忽略条件与结论的否定形式}：写否命题或逆否命题时，要正确写出 \(\lnot p\) 和 \(\lnot q\)，注意"大于"的否定是"不大于"（即小于等于），等等。
    \item \textbf{混淆"互逆"与"互否"关系}：可借助"交换"和"取反"的记忆：逆命题交换位置，否命题同时取反，逆否命题交换+取反。
\end{itemize}
\end{trapbox}